\section{Background}
\label{sec:background}

\subsection{Software Architecture Erosion}
Software architecture comprises the structures of a system, including
its elements, the relationships among them, and the properties of both
\cite{bass2012}. Perry and Wolf characterize architecture in terms of
elements, form, and rationale, emphasizing that architectural decisions
guide system development and should remain binding as the system evolves
\cite{perry1992}. Architecture erosion occurs when the implementation
progressively diverges from these intended structures and constraints.
Typical manifestations include layering violations, cyclic dependencies,
and unauthorized couplings
\cite{desilva2012,perry1992,passos2010}. Such deviations can undermine
maintainability and increase the cost of subsequent changes
\cite{bass2012,passos2010}.
\subsection{Architectural Fitness Functions and CI Enforcement}

Architectural fitness functions are executable checks that assess whether
selected architectural characteristics remain satisfied as a system
evolves \cite{ford2017}. Static conformance checks implement this idea by
evaluating source code or compiled artifacts against explicit
architectural constraints \cite{passos2010}. Integrated into continuous
integration (CI), these checks can detect violations automatically and
enforce architectural constraints as part of the build and test process
\cite{pinto2017}.

In this study, we examines enforcement through GitHub Actions workflows on the
default branch. A check acts as an integration gate when a workflow
invokes the relevant checker, test, or build task and its result can
cause the workflow to fail. Thus, defining an architectural rule does not
necessarily mean that it is enforced: enforcement depends on both the
rule's configuration and its execution within the CI workflow.

\subsection{Architecture Checkers and Enforcement Mechanisms}

Our study examines three architectural conformance tools:
ArchUnit for JVM-based systems, import-linter for Python, and
dependency-cruiser for JavaScript and TypeScript. Although all three
express constraints on software dependencies and structure, they differ
in how rules are defined, invoked, and selectively relaxed.

\paragraph{ArchUnit.}
ArchUnit expresses architectural constraints as tests over compiled
Java bytecode and JVM classes \cite{archunit}. Its rules can constrain
dependency directions, enforce layer relationships, detect cycles, and
restrict classes according to naming conventions, annotations, or
visibility. A rule violation causes the corresponding test to fail,
unless the rule's enforcement has been relaxed or disabled.

ArchUnit supports several mechanisms for handling violations.
A \texttt{FreezingArchRule} records existing violations as a baseline
and allows them to persist while detecting newly introduced violations.
The \texttt{@ArchIgnore} annotation can exclude a rule from execution.
In addition, rule severity can be configured so that violations no
longer cause test failure. ArchUnit does not require a dedicated
checker-specific CI command: it is typically invoked through a build
or test task, such as Maven's \texttt{test}, \texttt{verify},
\texttt{install}, or \texttt{package}, or Gradle's \texttt{test},
\texttt{check}, or \texttt{build}. Accordingly, CI integration is
identified through the execution of these build or test tasks.

\paragraph{import-linter.}
import-linter checks import relationships in Python projects
\cite{importlinter}. Its contracts can specify which modules may import
one another, constrain dependencies between layers, and prohibit
circular imports. The checker is typically invoked through the
\texttt{lint-imports} command, with its contracts defined in a
configuration file.

A contract can use an \texttt{ignore\_imports} list to exempt selected
import relationships from checking. This mechanism narrows the scope
of enforcement while retaining the contract and its remaining constraints
in the repository. In CI, import-linter is enforced when a workflow
step invokes the checker and propagates its result to the workflow.

\paragraph{dependency-cruiser.}
dependency-cruiser validates dependency relationships among JavaScript
and TypeScript modules \cite{depcruise}. Its rules can allow or forbid
particular dependencies and identify circular dependencies. Rules also
carry severity settings that determine how violations affect the
checker’s result.

Enforcement can be narrowed by excluding selected paths from a rule.
Alternatively, severity can be lowered so that violations are reported
without causing the check to fail, depending on the configured severity.
The tool is invoked through the \texttt{depcruise} command or an
equivalent project-specific command in a CI workflow. As with
import-linter, the configuration defines the constraints, whereas the
workflow establishes whether those constraints are checked during
integration.

\subsection{Enforcement Narrowing: Exemptions and Disabling}

Despite their different configuration mechanisms, these tools
provide two broad ways to weaken architectural enforcement:
\emph{exemptions} and \emph{disabling}. These categories describe the
effect of a configuration on enforcement rather than a common syntax
shared by the tools.

An \emph{exemption} retains the check but narrows the violations that
it treats as failures. In ArchUnit, a \texttt{FreezingArchRule} permits
violations already captured in a baseline while detecting new ones.
In import-linter, \texttt{ignore\_imports} excludes selected import
relationships. In dependency-cruiser, path-based exceptions exclude
selected modules from a rule. Lowering a rule's severity can also
preserve violation reporting while removing the failure condition.
Although these mechanisms differ in scope and behavior, they all allow
some violations to remain without necessarily failing the CI check.

\emph{Disabling}, by contrast, removes a rule or its violations from
the set of conditions that can cause failure. Examples include skipping
an ArchUnit rule with \texttt{@ArchIgnore}, configuring a rule so that
its violations no longer affect the outcome, or deleting or commenting
out the rule. The precise effect depends on the tool and configuration:
a disabled rule may no longer execute, whereas a rule with non-failing
severity may still execute and report violations.

The distinction between exemptions and disabling therefore depends on
what remains enforced after a configuration change. An exemption
preserves enforcement over a restricted scope, while disabling removes
the relevant rule or violation from failure-producing enforcement.
In both cases, the rule definition may remain in the repository even
when its practical enforcement has been weakened.

We define \emph{rule rot} as this weakening: an edit that exempts
or disables an architectural check, so the violations that can fail
a build become fewer, or the check is removed from failure-producing
enforcement. The rule specification may remain after the edit.

Finally, rule configuration and CI integration provide complementary
evidence of enforcement. The rule configuration identifies the
architectural constraint and any exemptions or disabling mechanisms,
whereas the workflow establishes whether the relevant check is invoked.
For ArchUnit, this requires tracing the workflow to a build or test task
that executes the architectural tests; for import-linter and
dependency-cruiser, it requires identifying the command that invokes
the checker. Considering both sources makes it possible to distinguish
the presence of architectural rules from their effective enforcement
during CI.